\chapter[Exact Ground-State Models]{Exact Ground-State Models: AKLT, Cluster, and Parent Hamiltonians}
\label{chap:exact-ground-state-models}

The preceding chapters presented several different notions of exact solvability
in two-dimensional quantum many-body systems.  The toric code and quantum double
models are commuting-projector models, so their full spectra reduce to local
constraint eigenvalues.  The Kitaev honeycomb model is not a commuting-projector
model, but it reduces to free Majorana fermions in each static gauge sector.
Rokhsar--Kivelson quantum dimer models provide exact ground states at special
frustration-free points, while the full spectrum is generally not solved.

This chapter develops a broader exact-ground-state viewpoint.  The central
objects are Hamiltonians whose ground states are known exactly because the local
terms are projectors annihilating a prescribed many-body state.  The paradigms
are
\begin{enumerate}[label=(\roman*)]
    \item the AKLT chain, whose valence-bond-solid ground state is an exact
    matrix-product state;
    \item the cluster model, whose ground state is a stabilizer state and a
    simple example of symmetry-protected topological order;
    \item the general parent-Hamiltonian construction for MPS and PEPS.
\end{enumerate}
The common mechanism is
\boxedpar{%
exactly known wavefunction
\(\Longrightarrow\)
local reduced-density-matrix support
\(\Longrightarrow\)
positive local projector
\(\Longrightarrow\)
frustration-free Hamiltonian.}
This is a different sense of solvability from Bethe ansatz or free-fermion
diagonalization.  It is exact and nonperturbative, but it usually controls the
ground-state subspace rather than every excited level.

\section{Frustration-free exact ground-state solvability}
\label{sec:frustration-free-ground-state-solvability}

Let a local Hamiltonian be written as
\begin{equation}
\label{eq:ff-general-hamiltonian}
    H=\sum_X h_X,
    \qquad
    h_X\geq 0,
\end{equation}
where each term acts on a finite region \(X\) of the lattice.  A state
\(\ket{\Psi}\) is a frustration-free ground state if
\begin{equation}
\label{eq:ff-ground-condition}
    h_X\ket{\Psi}=0
    \qquad
    \text{for every local term }X.
\end{equation}
Then
\begin{equation}
\label{eq:ff-zero-energy-ground-state}
    H\ket{\Psi}=0,
    \qquad
    H\geq0,
\end{equation}
so \(\ket{\Psi}\) is a ground state with energy zero.  If the kernel of
\(H\) is known exactly, the ground-state problem is solved.

\begin{definition}[Frustration-free Hamiltonian]
A positive local Hamiltonian \(H=\sum_X h_X\), with \(h_X\geq0\), is called
frustration-free if
\begin{equation}
\label{eq:ff-kernel-intersection}
    \ker H
    =
    \bigcap_X \ker h_X .
\end{equation}
Equivalently, a global ground state minimizes every local term separately.
\end{definition}

This definition should be compared with commuting-projector solvability.  If all
local projectors commute, then their common eigenspaces give the full spectrum.
If the terms are frustration-free but noncommuting, the ground space can still be
known exactly, but the excited spectrum is generally nontrivial.

\begin{remark}[Three strengths of solvability]
It is useful to distinguish three levels.
\begin{enumerate}[label=(\alph*)]
    \item \emph{Full-spectrum commuting-projector solvability}: all local terms
    commute, as in the toric code.
    \item \emph{Free-particle or Bethe-ansatz spectral solvability}: the full
    spectrum is generated by quasiparticles or rapidities, as in the XY and XXZ
    chains.
    \item \emph{Exact ground-state solvability}: the ground-state manifold is
    exactly known, but the full spectrum is not generally diagonalized.  AKLT,
    cluster, RK, and generic MPS/PEPS parent Hamiltonians mostly belong to this
    class.
\end{enumerate}
\end{remark}

The exact-ground-state construction is especially powerful because it produces
nontrivial phases with simple wavefunctions.  In one dimension, it naturally
leads to matrix-product states.  In two dimensions, it leads to projected
entangled-pair states, string-net states, dimer liquid states, and other tensor-
network fixed points.

\section{The AKLT chain}
\label{sec:aklt-chain}

The Affleck--Kennedy--Lieb--Tasaki (AKLT) chain is the canonical exactly soluble
spin-1 model.  It realizes the Haldane phase by an explicit valence-bond-solid
(VBS) wavefunction.  Its solvability is not based on free fermions or Bethe
rapidities.  Instead, it follows from a projector decomposition of two adjacent
spin-1 degrees of freedom.

\subsection{Spin-1 degrees of freedom and two-site projectors}
\label{subsec:aklt-two-site-projectors}

Let \(\bm S_j=(S_j^x,S_j^y,S_j^z)\) be spin-1 operators satisfying
\begin{equation}
\label{eq:spin-one-algebra}
    [S_j^\alpha,S_\ell^\beta]
    =
    \ii\delta_{j\ell}\epsilon^{\alpha\beta\gamma}S_j^\gamma,
    \qquad
    \bm S_j^2=2\id .
\end{equation}
For two neighboring spin-1 sites, the tensor product decomposes as
\begin{equation}
\label{eq:spin-one-two-site-decomposition}
    1\otimes1 = 0\oplus1\oplus2 .
\end{equation}
Let \(P_{j,j+1}^{(S)}\) denote the projector onto total spin \(S\) in the pair
\((j,j+1)\).  Since
\begin{equation}
\label{eq:spin-dot-eigenvalues-spin-one}
    \bm S_j\cdot\bm S_{j+1}
    =
    \frac12
    \left[
        (\bm S_j+\bm S_{j+1})^2
        -\bm S_j^2
        -\bm S_{j+1}^2
    \right],
\end{equation}
its eigenvalues in the \(S=0,1,2\) sectors are
\begin{equation}
\label{eq:spin-one-dot-eigenvalues}
    x_0=-2,
    \qquad
    x_1=-1,
    \qquad
    x_2=1 .
\end{equation}
Therefore the projector onto the total-spin-2 sector is the polynomial
\begin{equation}
\label{eq:spin-two-projector-polynomial}
    P_{j,j+1}^{(2)}
    =
    \frac{1}{6}
    \left(
        \bm S_j\cdot\bm S_{j+1}+2
    \right)
    \left(
        \bm S_j\cdot\bm S_{j+1}+1
    \right).
\end{equation}
Equivalently,
\begin{equation}
\label{eq:spin-two-projector-expanded}
    P_{j,j+1}^{(2)}
    =
    \frac16(\bm S_j\cdot\bm S_{j+1})^2
    +\frac12\bm S_j\cdot\bm S_{j+1}
    +\frac13 .
\end{equation}

The AKLT Hamiltonian on a periodic chain is
\begin{equation}
\label{eq:aklt-hamiltonian-projector}
    H_{\rm AKLT}
    =
    \sum_{j=1}^{L} P_{j,j+1}^{(2)},
    \qquad
    L+1\equiv1 .
\end{equation}
Up to an overall factor and additive constant, this is the more familiar
bilinear-biquadratic form
\begin{equation}
\label{eq:aklt-hamiltonian-bilinear-biquadratic}
    H_{\rm AKLT}'
    =
    \sum_{j=1}^{L}
    \left[
        \bm S_j\cdot\bm S_{j+1}
        +\frac13(\bm S_j\cdot\bm S_{j+1})^2
        +\frac23
    \right]
    =
    2H_{\rm AKLT} .
\end{equation}
The projector form \cref{eq:aklt-hamiltonian-projector} is the most transparent
for exact solvability.

\subsection{Valence-bond-solid construction}
\label{subsec:aklt-vbs-construction}

The VBS construction replaces every physical spin-1 by two virtual spin-
\(1/2\) degrees of freedom.  At site \(j\), the two virtual spins are projected
onto their symmetric triplet subspace:
\begin{equation}
\label{eq:spin-one-as-symmetric-two-spin-half}
    \calH^{S=1}_j
    =
    \operatorname{Sym}
    \left(
        \CC^2_{j,L}\otimes\CC^2_{j,R}
    \right).
\end{equation}
Neighboring virtual spins form singlets
\begin{equation}
\label{eq:virtual-singlet}
    \ket{\omega}_{j,R;j+1,L}
    =
    \frac{1}{\sqrt2}
    \left(
        \ket{\uparrow}_{j,R}\ket{\downarrow}_{j+1,L}
        -
        \ket{\downarrow}_{j,R}\ket{\uparrow}_{j+1,L}
    \right).
\end{equation}
The AKLT state is obtained by placing singlets on all virtual bonds and applying
the triplet projection at every site:
\begin{equation}
\label{eq:aklt-vbs-state}
    \ket{\Psi_{\rm AKLT}}
    =
    \prod_j P^{\rm sym}_j
    \prod_j
    \ket{\omega}_{j,R;j+1,L} .
\end{equation}
For an open chain, two unpaired virtual spin-\(1/2\) degrees of freedom remain
at the ends.  They produce a four-dimensional edge multiplet before boundary
terms are added.  This is the simplest microscopic origin of the spin-
\(1/2\) edge states in the Haldane phase.

\begin{proposition}[AKLT ground state]
The VBS state \(\ket{\Psi_{\rm AKLT}}\) is annihilated by every local projector
\(P_{j,j+1}^{(2)}\).  Hence it is a zero-energy ground state of
\(H_{\rm AKLT}\).
\end{proposition}

\begin{proof}
Consider two neighboring physical spin-1 sites.  In the virtual construction,
the right virtual spin of site \(j\) and the left virtual spin of site \(j+1\)
form a singlet.  Thus the total spin carried by the four virtual spin-
\(1/2\) variables on the two sites cannot reach total spin 2.  After projection
to the physical triplet subspace at each site, the two physical spin-1's
therefore have no component in the total-spin-2 sector.  Consequently
\(P_{j,j+1}^{(2)}\ket{\Psi_{\rm AKLT}}=0\) for every bond.  Since the Hamiltonian
is a sum of positive projectors, the state is a ground state.
\end{proof}

\subsection{AKLT state as a matrix-product state}
\label{subsec:aklt-mps}

The VBS state can be written as an MPS with bond dimension \(D=2\).  In the
basis \(\ket{+1},\ket{0},\ket{-1}\) of \(S^z\), one convenient convention is
\begin{equation}
\label{eq:aklt-mps-matrices}
    A^{+1}=\sqrt{\frac23}\,\sigma^+,
    \qquad
    A^{0}=-\frac{1}{\sqrt3}\,\sigma^z,
    \qquad
    A^{-1}=-\sqrt{\frac23}\,\sigma^- .
\end{equation}
The periodic-chain wavefunction is
\begin{equation}
\label{eq:aklt-mps-periodic}
    \ket{\Psi_{\rm AKLT}}
    =
    \sum_{m_1,\ldots,m_L}
    \Tr\left(A^{m_1}A^{m_2}\cdots A^{m_L}\right)
    \ket{m_1m_2\cdots m_L} .
\end{equation}
For an open chain, the trace is replaced by boundary spinors in the virtual
space.  The different boundary spinors describe the edge spin-
\(1/2\) degrees of freedom.

The MPS transfer channel is
\begin{equation}
\label{eq:aklt-transfer-channel}
    \mathbb{E}(X)
    =
    \sum_{m=-1}^{1} A^m X (A^m)^\dagger .
\end{equation}
With the convention \cref{eq:aklt-mps-matrices}, the identity is the leading
fixed point and the three Pauli matrices form subleading eigenoperators with
eigenvalue \(-1/3\).  Thus connected spin correlations decay as
\begin{equation}
\label{eq:aklt-correlation-length}
    \langle S_i^\alpha S_j^\alpha\rangle_c
    \propto
    \left(-\frac13\right)^{\abs{i-j}},
    \qquad
    \xi_{\rm AKLT}=\frac{1}{\log3} .
\end{equation}
This is an exact statement obtained directly from the finite-dimensional MPS
transfer matrix.

\begin{remark}[Hidden order]
Although ordinary spin correlations are short-ranged, the AKLT state has
nonzero string order.  A standard string correlator is
\begin{equation}
\label{eq:aklt-string-order}
    \mathcal{O}_{\rm string}^z
    =
    -\lim_{\abs{i-j}\to\infty}
    \left\langle
        S_i^z
        \exp\left(\ii\pi\sum_{k=i+1}^{j-1}S_k^z\right)
        S_j^z
    \right\rangle .
\end{equation}
For the AKLT state this quantity is finite.  It detects the hidden antiferromag-
netic order associated with the Haldane symmetry-protected topological phase.
\end{remark}

\section{The cluster model}
\label{sec:cluster-model}

The cluster model is the simplest stabilizer example of an exact ground-state
model with symmetry-protected topological structure.  Unlike the AKLT chain,
its local terms commute, so it is also full-spectrum soluble as a stabilizer
Hamiltonian.  It is nevertheless conceptually useful in this chapter because its
exact ground state is a graph state and therefore a tensor-network fixed point.

\subsection{One-dimensional cluster Hamiltonian}
\label{subsec:one-dimensional-cluster-hamiltonian}

Consider a chain of qubits with Pauli operators \(X_j,Y_j,Z_j\).  On a periodic
chain, define the stabilizers
\begin{equation}
\label{eq:cluster-stabilizers-periodic}
    K_j
    =
    Z_{j-1}X_jZ_{j+1},
    \qquad
    j=1,\ldots,L,
\end{equation}
with site labels understood modulo \(L\).  The one-dimensional cluster
Hamiltonian is
\begin{equation}
\label{eq:cluster-hamiltonian-periodic}
    H_{\rm cl}
    =
    -\sum_{j=1}^{L}K_j .
\end{equation}
The stabilizers commute:
\begin{equation}
\label{eq:cluster-stabilizer-commutation}
    [K_i,K_j]=0,
    \qquad
    K_j^2=\id .
\end{equation}
Indeed, if \(K_i\) and \(K_j\) overlap nontrivially, they either share only
commuting Pauli matrices or have two anticommuting overlaps, whose signs cancel.
Therefore the full spectrum is labelled by eigenvalues \(k_j=\pm1\).

The cluster ground state \(\ket{C}\) is the unique common \(+1\) eigenstate on a
periodic chain:
\begin{equation}
\label{eq:cluster-ground-state-stabilizer}
    K_j\ket{C}=\ket{C},
    \qquad
    j=1,\ldots,L .
\end{equation}
Equivalently, it is obtained by applying controlled-\(Z\) gates to the product
state \(\ket{+}^{\otimes L}\):
\begin{equation}
\label{eq:cluster-state-cz-construction}
    \ket{C}
    =
    \prod_{\langle j,j+1\rangle} CZ_{j,j+1}\,
    \ket{+}^{\otimes L},
    \qquad
    X_j\ket{+}_j=\ket{+}_j .
\end{equation}
This formula follows because conjugation by \(CZ_{j,j+1}\) maps
\begin{equation}
\label{eq:cz-pauli-conjugation}
    X_j
    \longmapsto
    Z_{j-1}X_jZ_{j+1}
\end{equation}
when all nearest-neighbor controlled-\(Z\) gates are applied.

\begin{remark}[Open-chain edge degeneracy]
For an open chain, one may define bulk stabilizers
\(K_j=Z_{j-1}X_jZ_{j+1}\) for \(2\leq j\leq L-1\).  Then the two boundary qubits
are not fully fixed by the bulk terms, and protected edge degrees of freedom
appear.  This parallels the AKLT edge spin-
\(1/2\)'s, although the cluster model is built from qubits and commuting
stabilizers rather than spin-1 projectors.
\end{remark}

\subsection{Symmetry-protected topological interpretation}
\label{subsec:cluster-spt}

The one-dimensional cluster state is a fixed point of a nontrivial
\(\ZZ_2\times\ZZ_2\) symmetry-protected topological phase.  A convenient choice
of symmetry generators is
\begin{equation}
\label{eq:cluster-z2xz2-symmetry}
    G_{\rm even}=\prod_{j\,\text{even}}X_j,
    \qquad
    G_{\rm odd}=\prod_{j\,\text{odd}}X_j .
\end{equation}
On a closed chain, these commute with the Hamiltonian.  On an open chain, the
symmetry acts projectively on the edge degrees of freedom.  This projective edge
action is the one-dimensional SPT invariant.

The cluster state also has a string order parameter.  For example, products of
stabilizers imply that
\begin{equation}
\label{eq:cluster-string-order}
    Z_i
    \left(\prod_{k=i+1}^{j-1}X_k\right)
    Z_j
\end{equation}
has long-range expectation value in the cluster ground state for suitable parity
choices of \(i\) and \(j\).  The precise endpoint operators depend on whether
one chooses the even or odd symmetry sector, but the qualitative structure is
robust: nonlocal string order replaces ordinary local order.

\subsection{Graph-state generalization}
\label{subsec:graph-state-generalization}

The cluster construction extends immediately to any graph \(\Gamma=(V,E)\).
Place a qubit on every vertex and define
\begin{equation}
\label{eq:graph-state-stabilizers}
    K_v
    =
    X_v\prod_{u:(u,v)\in E}Z_u .
\end{equation}
The graph-state Hamiltonian is
\begin{equation}
\label{eq:graph-state-hamiltonian}
    H_\Gamma=-\sum_{v\in V}K_v .
\end{equation}
All \(K_v\)'s commute, and the graph state is the common \(+1\) eigenstate.
For a two-dimensional square-lattice graph, this gives the two-dimensional
cluster state, an important stabilizer resource state for measurement-based
quantum computation.

\begin{remark}[Cluster model versus toric code]
Both the cluster model and the toric code are stabilizer Hamiltonians, but their
orders are different.  The cluster state is short-range entangled and becomes a
trivial product state after a finite-depth circuit if the protecting symmetry is
ignored.  The toric code is long-range entangled and has intrinsic topological
order.  Thus stabilizer solvability by itself does not determine the phase;
one must also examine entanglement, symmetry, and nonlocal operators.
\end{remark}

\section{Parent Hamiltonians of matrix-product states}
\label{sec:mps-parent-hamiltonians}

The AKLT chain is not an isolated miracle.  It is an example of a general MPS
parent-Hamiltonian construction.  The input is an exact tensor-network state;
the output is a local frustration-free Hamiltonian for which that tensor network
is an exact ground state.

\subsection{Matrix-product states}
\label{subsec:mps-definition-parent}

Let each site have physical dimension \(d\), with basis
\(\ket{s}\), \(s=1,\ldots,d\).  A translation-invariant MPS with bond dimension
\(D\) on a periodic chain is
\begin{equation}
\label{eq:mps-periodic-definition}
    \ket{\Psi_A}
    =
    \sum_{s_1,\ldots,s_L=1}^{d}
    \Tr\left(A^{s_1}A^{s_2}\cdots A^{s_L}\right)
    \ket{s_1s_2\cdots s_L},
\end{equation}
where each \(A^s\) is a \(D\times D\) matrix.  The transfer channel is
\begin{equation}
\label{eq:mps-transfer-channel-general}
    \mathbb{E}(X)
    =
    \sum_{s=1}^{d}A^sX(A^s)^\dagger .
\end{equation}
The state is called injective, possibly after blocking a finite number of sites,
if the blocked tensors generate the full matrix algebra on the virtual space.
More explicitly, for some finite length \(\ell\), the map
\begin{equation}
\label{eq:mps-injectivity-map}
    \Gamma_\ell:
    \operatorname{Mat}_{D\times D}(\CC)
    \to
    (\CC^d)^{\otimes \ell},
    \qquad
    X
    \mapsto
    \sum_{s_1,\ldots,s_\ell}
    \Tr\left(XA^{s_1}\cdots A^{s_\ell}\right)
    \ket{s_1\cdots s_\ell}
\end{equation}
is injective.

Injectivity is the tensor-network analogue of having a unique locally generated
ground state.  Noninjective MPS describe symmetry-breaking or topologically
ordered fixed-point structures and require a slightly refined parent
construction.

\subsection{Local support and the parent projector}
\label{subsec:mps-parent-projector}

Choose a block length \(\ell\) at least as large as the injectivity length.  Let
\(\mathcal{S}_\ell\) be the support subspace of the \(\ell\)-site reduced density
matrix of \(\ket{\Psi_A}\).  Equivalently,
\begin{equation}
\label{eq:mps-local-support-subspace}
    \mathcal{S}_\ell
    =
    \operatorname{Im}\Gamma_\ell .
\end{equation}
Let \(\Pi_\ell\) be the orthogonal projector onto \(\mathcal{S}_\ell\).  The local
parent-Hamiltonian term is
\begin{equation}
\label{eq:mps-parent-local-term}
    h_{j,j+\ell-1}
    =
    \id-\Pi_{j,j+\ell-1},
\end{equation}
where \(\Pi_{j,j+\ell-1}\) is the translated projector onto the corresponding
support on sites \(j,\ldots,j+\ell-1\).  The parent Hamiltonian is
\begin{equation}
\label{eq:mps-parent-hamiltonian}
    H_{\rm par}
    =
    \sum_j h_{j,j+\ell-1} .
\end{equation}
Because every local reduced density matrix of \(\ket{\Psi_A}\) is supported in
\(\mathcal{S}_\ell\), one has
\begin{equation}
\label{eq:mps-parent-annihilation}
    h_{j,j+\ell-1}\ket{\Psi_A}=0
    \qquad
    \text{for all }j .
\end{equation}
Thus \(\ket{\Psi_A}\) is an exact frustration-free ground state.

\begin{theorem}[Parent Hamiltonian of an injective MPS]
For an injective translation-invariant MPS, the parent Hamiltonian
\cref{eq:mps-parent-hamiltonian} has the MPS as its unique ground state on a
periodic chain, for sufficiently large system size.  On an open chain, the
ground-state space is controlled by the virtual boundary degrees of freedom.
\end{theorem}

\begin{remark}[AKLT as an MPS parent Hamiltonian]
For the AKLT MPS, the two-site physical Hilbert space decomposes into total spin
\(S=0,1,2\).  The two-site support of the VBS state contains only \(S=0\) and
\(S=1\), while the \(S=2\) subspace is absent.  Therefore the parent local term
is precisely the spin-2 projector \(P_{j,j+1}^{(2)}\).  The AKLT Hamiltonian is
an MPS parent Hamiltonian.
\end{remark}

\subsection{Spectral gap and correlation length}
\label{subsec:mps-gap-correlation}

For injective MPS in one dimension, the parent Hamiltonian is generically
expected to be gapped, and the standard injective parent construction indeed
produces a finite spectral gap in the thermodynamic limit.  The finite bond
dimension also implies exponential decay of connected correlations.  Both
features are controlled by the transfer channel \(\mathbb{E}\): if the leading
eigenvalue is normalized to one and the subleading eigenvalue has magnitude
\(\abs{\lambda_2}<1\), then typical connected correlators decay as
\begin{equation}
\label{eq:mps-correlation-length-general}
    C(r)
    \sim
    \lambda_2^r,
    \qquad
    \xi=-\frac{1}{\log\abs{\lambda_2}} .
\end{equation}
The parent Hamiltonian and the transfer channel are therefore two sides of the
same exact-ground-state structure: the parent Hamiltonian fixes the state as a
ground state, while the transfer channel computes its equal-time observables.

\section{Parent Hamiltonians of PEPS}
\label{sec:peps-parent-hamiltonians}

Projected entangled-pair states are the natural two-dimensional generalization
of MPS.  A PEPS assigns a tensor to each vertex of a lattice.  Its physical
index labels the local Hilbert space, while its virtual indices are contracted
along edges.  The AKLT construction itself has a two-dimensional analogue on
suitable lattices: place virtual spin-
\(1/2\)'s on incident bonds, form virtual singlets on edges, and project the
virtual degrees of freedom at each site into a physical spin representation.

The PEPS parent-Hamiltonian construction parallels the MPS case.  For a finite
region \(R\), let \(\mathcal{S}_R\) be the support of the reduced density matrix
of the PEPS on \(R\).  Define a local term
\begin{equation}
\label{eq:peps-parent-local-term}
    h_R=\id-\Pi_R,
\end{equation}
where \(\Pi_R\) projects onto \(\mathcal{S}_R\).  Then
\begin{equation}
\label{eq:peps-parent-hamiltonian}
    H_{\rm PEPS}=
    \sum_R h_R
\end{equation}
is frustration-free and has the PEPS as a ground state.

In two dimensions, however, the physics is more subtle than in injective
one-dimensional MPS.  PEPS parent Hamiltonians can describe
\begin{enumerate}[label=(\roman*)]
    \item short-range-entangled trivial phases;
    \item symmetry-protected topological phases;
    \item intrinsic topological order, when the tensors have virtual gauge
    symmetries;
    \item critical states, where correlations may decay algebraically;
    \item gapless or gapped parent Hamiltonians, depending on the tensor.
\end{enumerate}
Thus the PEPS parent construction gives an exact ground state, but it does not
by itself guarantee a simple spectrum or even a finite gap.

\begin{remark}[Virtual symmetries and topological order]
In PEPS, intrinsic topological order is often encoded by symmetries acting on
the virtual indices.  The toric code, string-net fixed points, and many quantum
double states admit PEPS descriptions with such virtual gauge structures.  In
this language, topological degeneracy arises from nontrivial virtual boundary
conditions rather than from ordinary local order parameters.
\end{remark}

\section{Comparison of AKLT, cluster, RK, and toric-code solvability}
\label{sec:exact-ground-state-comparison}

The models in this part use related but distinct exact mechanisms.  The
following table summarizes the differences.

\begin{center}
\begin{adjustbox}{max width=\textwidth}
\begin{tabular}{L{0.21\textwidth}L{0.25\textwidth}L{0.24\textwidth}L{0.22\textwidth}}
\toprule
Model or construction & Exact mechanism & What is exactly controlled & Full spectrum? \\
\midrule
Toric code & Commuting stabilizer projectors & Ground space, anyon sectors, all stabilizer energies & Yes, in stabilizer variables \\
Cluster model & Commuting graph-state stabilizers & Unique or edge-degenerate stabilizer ground state, SPT edge structure & Yes, in stabilizer variables \\
AKLT chain & Noncommuting projector parent Hamiltonian & VBS/MPS ground state, edge spin-
\(1/2\)'s, correlations & Not by the parent construction alone \\
RK dimer model & Frustration-free projector at RK point & Equal-amplitude or weighted dimer ground state, classical correlators & Generally no \\
Injective MPS parent Hamiltonian & Local support projector & Unique exact MPS ground state and exponential correlations & Generally no \\
PEPS parent Hamiltonian & Local support projector with possible virtual symmetry & Exact tensor-network ground space, possibly topological sectors & Generally no \\
\bottomrule
\end{tabular}
\end{adjustbox}
\end{center}

The important lesson is that projector methods are broader than commuting-
projector methods.  Commutativity solves the spectrum.  Frustration-free parent
Hamiltonians solve the ground state.  Tensor-network transfer matrices then
compute equal-time observables, much as the row-to-row transfer matrix computed
classical correlations in earlier chapters.

\section{A unified parent-Hamiltonian principle}
\label{sec:unified-parent-hamiltonian-principle}

The parent-Hamiltonian construction can be stated abstractly.

\begin{principle}[Parent-Hamiltonian principle]
Given a many-body state \(\ket{\Psi}\), choose local regions \(R\) and let
\(\Pi_R\) project onto the support of the reduced density matrix
\(\rho_R=\Tr_{\overline R}\ket{\Psi}\bra{\Psi}\).  Then
\begin{equation}
\label{eq:abstract-parent-hamiltonian}
    H_{\rm par}[\Psi]
    =
    \sum_R(\id-\Pi_R)
\end{equation}
is positive semidefinite and satisfies
\begin{equation}
\label{eq:abstract-parent-annihilation}
    H_{\rm par}[\Psi]\ket{\Psi}=0 .
\end{equation}
If the intersection of the local supports is exactly the desired ground-state
space, then \(H_{\rm par}[\Psi]\) is an exact frustration-free Hamiltonian for
that state or phase.
\end{principle}

This principle explains why the AKLT chain, graph states, PEPS string-net fixed
points, and many RK-like constructions fit naturally into the same lecture-note
architecture.  They are not solved by the same algebra, but they share the same
logic:
\boxedpar{%
local tensor or local constraint
\(\Longrightarrow\)
finite-dimensional support condition
\(\Longrightarrow\)
projector Hamiltonian
\(\Longrightarrow\)
exact ground state and controlled correlations.}

\section{Common pitfalls}
\label{sec:aklt-cluster-parent-pitfalls}

Several points are worth keeping separate.

\begin{enumerate}[label=(\roman*)]
    \item \emph{Frustration-free does not mean commuting.}  The AKLT local
    projectors do not commute on overlapping bonds, although the exact ground
    state is annihilated by all of them.
    \item \emph{Exact wavefunction does not imply exact spectrum.}  Parent
    Hamiltonians usually determine the ground-state subspace, not all excited
    eigenstates.
    \item \emph{Short-range entangled does not mean trivial under symmetry.}
    AKLT and cluster states are short-range entangled, but they represent
    nontrivial SPT phases when the relevant symmetries are preserved.
    \item \emph{Tensor-network exactness is algebraic, not variational.}  In an
    exact parent-Hamiltonian construction, the tensor network is not merely a
    good approximation; it is annihilated exactly by local projectors.
    \item \emph{Two-dimensional parent Hamiltonians are more delicate than
    one-dimensional ones.}  PEPS may have topological sectors, virtual gauge
    symmetries, critical correlations, or gapless parent Hamiltonians.
\end{enumerate}

\section{Connection to the rest of the notes}
\label{sec:exact-ground-state-models-connections}

The AKLT, cluster, and parent-Hamiltonian constructions complete the set of
solvability mechanisms emphasized in these notes.  They should be read in
parallel with the previous chapters as follows:
\begin{equation}
\label{eq:solvability-mechanism-chain-final}
\begin{array}{ccl}
\text{XY/Kitaev/SSH} &:& \text{quadratic fermions and BdG diagonalization},\\[2pt]
\text{2D Ising} &:& \text{transfer matrices and free-fermion structure},\\[2pt]
\text{XXZ/six-vertex} &:& \text{Yang--Baxter algebra and Bethe ansatz},\\[2pt]
\text{toric code/quantum double} &:& \text{commuting projectors and topological constraints},\\[2pt]
\text{Kitaev honeycomb} &:& \text{Majorana fermions in static gauge sectors},\\[2pt]
\text{RK/AKLT/cluster/parent models} &:& \text{exact frustration-free ground states}.
\end{array}
\end{equation}
Together these examples show that exact solvability is not a single technique.
It is a family of algebraic reductions, each exposing a different kind of
structure in many-body quantum systems.
